\fsection{Empresas del sector en el entorno}

Motilla del Palancar, con una población de aproximadamente 6.000 habitantes, se sitúa en la comarca de La Manchuela conquense, junto a la autovía A-3 (Madrid-Valencia). Su tejido industrial está orientado principalmente al sector agroalimentario, el transporte por carretera y los servicios. No existen en el municipio empresas especializadas en automatización industrial, aunque sí se identifican empresas con necesidades potenciales de automatización:

\begin{itemize}
	\item Extramoti S.L. (Motilla del Palancar): empresa dedicada a la fabricación, reparación y comercialización de maquinaria industrial, así como servicios de granallado, fresado y pulido. Representa un perfil de cliente potencial para soluciones de automatización de procesos de mecanizado.

	\item Logidos (Motilla del Palancar): empresa de servicios logísticos con cobertura nacional. Sector con alta demanda de automatización de picking y clasificación de paquetes.

	\item Cooperativa Agraria Inmaculada Concepción (Motilla del Palancar): cooperativa del sector agroalimentario con necesidades de manipulación y clasificación de producto.

	\item Congelados Margabal / Congelados Júcar (Motilla del Palancar): empresas del sector alimentario con líneas de distribución que pueden beneficiarse de sistemas de identificación y clasificación automatizada de producto.
\end{itemize}

En un radio más amplio, el corredor de la A-3 entre Motilla del Palancar y Valencia concentra un número significativo de parques logísticos y plantas industriales de los sectores del automóvil, alimentación y distribución, donde la automatización del picking y la trazabilidad de producto mediante RFID son demandas crecientes.
